\chapter{Elements of Concurrent Program Design}\label{chap:ed_elements_of_concurrent_program_design} % (fold)
Designing \textbf{concurrent programs} consists of two macro steps:
\begin{itemize}
    \item \textbf{step 1 - analysis}: identifying tasks;
    \item \textbf{step 2 - design}: concurrent architectures;
\end{itemize}
%
\begin{definition}{Analysis}{cp-analysis}
   Analyzing the problem to identify exploitable concurrency: identifying \textbf{task \& data} decomposition and their \textbf{dependencies}.
    \begin{definition}{Task}{task}
        A sequence of instructions that \red{operate together as a group}, corresponding to some logical part of an algorithm or a program.
        Also, it can be considered as a specific piece of work required to be done as a duty or responsibility of some agent: more abstract/design-oriented definition.
    \end{definition}
\end{definition}
\begin{definition}{Design}{cp-design}
    Choosing a suitable \textbf{concurrent architecture}, assigning tasks to ``agents'':
    \begin{itemize}
        \item \textit{active components} responsible for performing \hyperref[def:task]{Tasks};
        \item encapsulate login \& control of \hyperref[def:task]{Task} execution;
    \end{itemize}
    In order to work together, agents share and use some \textit{passive components}:
    \begin{itemize}
        \item \textit{shared resources} encapsulating result of agent work;
        \item \textit{coordination media/resources} to shape agent interaction and manage \hyperref[def:task]{Task} dependencies;
    \end{itemize}
\end{definition}
\begin{note}[Active vs Passive Component Discipline]
    Strong partition:
    \begin{itemize}
        \item \textbf{active components} (``agents''):
            \begin{itemize}
                \item encapsulating a thread of control;
                \item no public interface;
                \item every interaction enabled and mediated by shared passive components;
            \end{itemize}
        \item \textbf{passive components}, designed as \hyperref[def:monitor]{Monitors};
    \end{itemize}
    \begin{warning}[No hybrid entities]
        An \textbf{hybrid entity} can be designed as a set of one or more active and passive components
    \end{warning}
\end{note}
%
\section{Problem Analysis}\label{sec:ed_problem_analysis} % (fold)
\begin{note}[Objective]
    Identifying \hyperref[def:task]{Task}/data decomposition and their dependencies
\end{note}
Usually performed in two steps:
\begin{itemize}
    \item \hyperref[def:task]{Task} and \textbf{data decomposition} analysis;
    \item Dependency Analysis;
\end{itemize}
%
\subsection{Task Decomposition Pattern}\label{sub:ed_task_decomposition_pattern} % (fold)
Following a \textit{divide-et-impera} principle, the problem can be naturally decomposed into a collection of independent or nearly independent \hyperref[def:task]{Tasks}.
\begin{note}
    \hyperref[def:task]{Tasks} must be sufficiently independent so as that managing dependencies takes only a small fraction of the program's overall execution time
\end{note}
First \hyperref[def:task]{Tasks} are identified, then how data can be decomposed given the task decomposition: in this case, \hyperref[sub:ed_data_decomposition_pattern]{Data Decomposition} follows \hyperref[sub:ed_task_decomposition_pattern]{Task Decomposition}.
\begin{example}{}{}
    \hyperref[sub:cp_producer_consumer_p_c_problem]{Producer-Consumer} systems.
\end{example}
% subsection Task Decomposition Pattern (end)
%
\subsection{Data Decomposition Pattern}\label{sub:ed_data_decomposition_pattern} % (fold)
Problem's data can be easily decomposed into units that can be operated on relatively independently.
First data units are identified, then \hyperref[def:task]{Task} related to that units can be identified: in this case, \hyperref[sub:ed_task_decomposition_pattern]{Task Decomposition} follows \textbf{data decomposition}.
\begin{note}
    Scale well with the available number of processors
\end{note}
\begin{example}{}{}
    Matrix Multiplication.
\end{example}
% subsection Data Decomposition Pattern (end)
%
\subsection{Dependency Analysis}\label{sub:ed_dependency_analysis} % (fold)
Analyzing dependencies among the \hyperref[def:task]{Tasks} within a problem's decomposition to \textbf{group} and \textbf{order} \hyperref[def:task]{Tasks} according to the type of dependencies, in order to maximize \textbf{performance}.
Two possible kind of dependencies:
\begin{itemize}
    \item \textbf{temporal} dependencies;
    \item \textbf{data/resource} dependencies;
\end{itemize}
% subsection Dependency Analysis (end)
% section Problem Analysis (end)
%
\section{Design Level}\label{sec:ed_design_level} % (fold)
From analysis to design:
\begin{itemize}
    \item \textbf{result parallelism}: envisioning parallelism in terms of program's result;
    \item \textbf{specialist parallelism}: envisioning parallelism in terms of an ensemble of \textit{specialists} that collectively constitute the programs;
    \item \textbf{agenda parallelism}: envisioning parallelism in terms of program's \textit{agenda} of \hyperref[def:task]{Tasks};
\end{itemize}
%
\subsection{Conceptual Classes}\label{sub:ed_conceptual_classes} % (fold)
\subsubsection{Result Parallelism}\label{sub:ed_result_parallelism} % (fold)
Designing the system \textit{around the data structure or resource yielded} as the ultimate result: we get parallelism by computing all the elements of the result simultaneously.
\begin{note}
    Strictly related to the \hyperref[sub:ed_data_decomposition_pattern]{Data Decomposition Pattern}
\end{note}
Each agent is assigned to produce one piece of the result: they all work in parallel up to the natural restriction imposed by the problem.
Proper \textbf{shared data structures} are designed to wrap the result under construction, encapsulating mutex and synchronization issues.
\begin{example}{Result Parallelism: Building a House}{}
    Think about the components of the house: front, rear, side walls, interior walls, foundations, roof, etc.
    Proceed by building all the components simultaneously, assembling them as they are completed.
    \emptyln{}
    Using separate agents (\textit{workers}) on the different parts, proceeding in parallel up to the point where work on one component can't proceed until another is finished.
\end{example}
% subsubsection Result Parallelism (end)
%
\subsubsection{Specialist Parallelism}\label{sub:ed_specialist_parallelism} % (fold)
Designing the system \textit{around an ensemble of specialists} connected into a logical network of some kind, parallelism results from all the \textit{nodes} of the logical network being active simultaneously.
\emptyln{}
Each agent is assigned to perform one specified kind of \hyperref[def:task]{Task}: they all work in parallel up to the natural restriction imposed by the problem.
\\
Coordination means (communication protocols, shared data structures) are used to support specialists communication and coordination.
\begin{note}
    Opposite approach with respect to result parallelism
\end{note}
\begin{example}{Specialist Parallelism: Building a House}{}
    Identify the separate skills required to build the house: surveyors, excavators, foundation builders, carpenters, roofers, etc.
    \begin{note}[Different kind of tasks]
    \end{note}
    Then assemble a constructor \textit{crew} in which each skill is represented by a separated specialist worker.
    \emptyln{}
    They all start simultaneously, but initially most \textbf{workers} will have to wait around.
    Only when the project is well underway, many \textit{workers} will be called to play simultaneously.
    \begin{warning}
        Strong role of worker interaction and coordination: different kind of strategies (e.g. pipelined jobs)
    \end{warning}
\end{example}
% subsubsection Specialist Parallelism (end)
%
\subsubsection{Agenda Parallelism}\label{sub:ed_agenda_parallelism} % (fold)
Designing the system around a particular \textbf{agenda of \hyperref[def:task]{Tasks}} and then assign many workers to each step, and also possibly carry on different steps in parallel (\textit{workflow}).
Each agent is assigned to help out with the current item on the agenda, by working in parallel up to the natural restriction imposed by the problem.
\emptyln{}
Shared data structures are designed to manage data consumed and produced by agenda activities (e.g. \textit{bounded-buffers}).
\begin{example}{Agenda Parallelism: Building a House}{}
    We write out a sequential agenda and carry it in order, but at each stage we assign mny workers to the current activity: doing the foundation, doing the frame, doing the roof, etc.
    \\
    We assemble a work team of \textbf{generalists}, each one capable of performing any construction step.
    \begin{note}
        The team is set to work stage by stage following the agenda
    \end{note}
\end{example}
% subsubsection Agenda Parallelism (end)
%
\paragraph{Remarks}\label{par:ed_cc_remarks} % (fold)
The boundaries between the three classes are not rigid, we will often mix elements of several approaches in getting a particular job done.
\begin{example}{}{}
    A \hyperref[sub:ed_specialist_parallelism]{Specialist Parallelism} approach might take secondary use of \hyperref[sub:ed_agenda_parallelism]{Agenda Parallelism}, for example by assigning a team of workers to some specialty.
\end{example}
However, the three approaches represent \textit{three clearly} separate ways of thinking about the problem:
\begin{itemize}
    \item \hyperref[sub:ed_result_parallelism]{Result Parallelism} focuses on the shape of the finished product;
    \item \hyperref[sub:ed_specialist_parallelism]{Specialist Parallelism} focuses on the makeup of the work crew;
    \item \hyperref[sub:ed_agenda_parallelism]{Agenda Parallelism} focuses on the list of \hyperref[def:task]{Tasks} to be performed;
\end{itemize}
\begin{note}
    The approach can be recursively applied, following \hyperref[sub:ed_task_decomposition_pattern]{Task Decomposition} and \hyperref[sub:ed_data_decomposition_pattern]{Data Decomposition} patterns
\end{note}
% paragraph Remarks (end)
%
\subsection{Using Conceptual Classes}\label{sub:ed_using_conceptual_classes} % (fold)
To write a \hyperref[def:conc-program]{Concurrent Program}:
\begin{itemize}
    \item choose the \hyperref[sub:ed_conceptual_classes]{Conceptual Class} that is most natural for the problem;
    \item write a program using the method that is most natural for that \hyperref[sub:ed_conceptual_classes]{Conceptual Class};
    \item if the resulting program is not acceptably efficient, transform it methodically into a more efficient version by switching from a more-natural method to a more-efficient one;
\end{itemize}
% subsection Using Conceptual Classes (end)
% subsection Conceptual Classes (end)
%
\subsection{Concurrent Architectures}\label{sub:ed_concurrent_architectures} % (fold)
Some specific architectures can help to bridge the gap between the \hyperref[sub:ed_conceptual_classes]{Conceptual Classes} and the implementation:
\begin{itemize}
    \item \textbf{Master-Workers}: \hyperref[def:task]{Task}-oriented frameworks, fork-join variants;
    \item \textbf{Filter-Pipeline}: stream-based patterns;
    \item \textbf{Shared-Space/Blackboard}: space-based coordination patterns;
    \item \textbf{Announcer-Listeners}: event-based coordination patterns;
\end{itemize}
%
\subsubsection{Master-Workers}\label{sub:ed_master_workers} % (fold)

% subsubsection Master-Workers (end)
% subsection Concurrent Architectures (end)
% section Design Level (end)
% chapter Elements of Concurrent Program Design (end)
